\scanpage{22}
% @ 1–3, top
Faisant des extensions de la base successives $Y_1=X$, [d'où $X_1^{\deleted{\unsure{*}}}=X\times_Y Y_1$, $X_1=X_{(1)}\amalg X^{(1)}$, puis $Y_2=X^{(1)}$,] on trouve

% @ 4–9, bracketed lemma
\uline{Lemme} 2 Il existe un $Y'\deleted{_{\unsure{2}}}$ \deleted{\unsure{sur}} fini sur $Y$, avec $Y'\deleted{_{\unsure{2}}}\to Y$ surjectif [et \deleted{radiciel} \inserted{univ\textsuperscript{t} ouvert} si on veut] tel que $X'=X\times_Y Y'$ se décompose en $n$ \unsure{variétés|revêtements|réunion} $X'_i$ connexes de degré \deleted{fini} séparable $1$ sur $Y$.

% @ 10–16, below the fold
Soit $I$ l'ens des composantes connexes de $X'$, \unsure{considérons|considérer} $Y'\times I$ comme rev. étale de $Y'$. Il est trivial que $X'_0=Y'\times I$ est muni d'une donnée de descente \unsure{relat\textsuperscript{vt} à|relativement à} $Y'\to Y$, d'où un \uline{rev.} \uline{étale} $X$ de $Y$, donnant $X'_0{}^{\deleted{\unsure{*}}}$ par image inverse.

% @ lower left, level with 15–18
\margin{\diagram{\begin{tikzcd}[ampersand replacement=\&,row sep=small,column sep=small]
X_{0}^{\deleted{\unsure{*}}}\arrow[d]\&\&X'_{0}{}^{\deleted{\unsure{*}}}\arrow[ll]\arrow[d]\\
Y\&\&Y'\arrow[ll]\\
\&Y^{*}_{\deleted{\unsure{2}}}\arrow[ul]\&\&Y'{}^{*}_{\deleted{\unsure{2}}}\arrow[ul]\arrow[ll]
\end{tikzcd}}}

% @ 17–19, foot
Soit ${\unsure{Y|Y'}}^{*}_{\deleted{\unsure{2}}}$ un rev. \uline{étale} [donc plat] de $Y$ qui décompose $X_0$ complètement, prenons l'image inverse de la situation sur
% @ sideways band through the manuscript
% printed form
\textsc{Extra Strong}\\
\unsure{Guerimand|Guérimand} Voiron
% archive sheet number: 21
